\subsection*{MT-ART09. A smaller relative-zero gate for solvable Artin holomorphy}

Let $L/K$ be a finite solvable Galois extension, with group $G$,
and let $s_0\ne1$. Write
\[
 q(s_0)=\operatorname{ord}_{s=s_0}
             \frac{\zeta_L(s)}{\zeta_K(s)}\ge0.
\]
The nonnegativity follows from the classical Aramata--Brauer
theorem. When $|G|$ has at least two distinct prime divisors,
let $p_2$ be the second smallest and put
\[
                 B(G)=\max\{2,p_2-2\}.
\]
Then every Artin $L(s,\chi,L/K)$ is holomorphic at $s_0$ if
$q(s_0)\le B(G)$. This is the combination of Peng-Jie Wong,
\emph{Zeros of Dedekind zeta functions and holomorphy of Artin
$L$-functions}, Theorems1.3 and1.4, printed p.254, with their
complete proofs on pp.259 and260. The solvability hypothesis
is essential to the statements as cited here. In a $p$-group
all irreducible characters are monomial, so Artin holomorphy
already follows globally from Hecke holomorphy.

For $4\nmid |G|$ the sufficient bound $q(s_0)\le3$ holds,
without adding a solvability assumption: Wong Corollary5.3,
printed pp.259--260, includes the group-theoretic implication
that these groups are solvable. His Corollary5.1, p.259, is the
odd-degree special case. One may therefore replace $B(G)$ by
$\max\{B(G),3\}$ when applicable. The stronger Sylow-structure
variants in his Theorems5.2 and5.4 are separately located there;
no unstated prime-factor hypothesis from them is being used.

\paragraph{Complete transfer to the adopted Hecke hypothesis.}
Under $M_H$, a possible pole of an Artin function in
$0<\Re s<1$ must lie on $\Re s=1/2$: Brauer induction writes
that function as a finite quotient of Hecke functions, whose
possible nontrivial denominator zeros are all on that line.
This is the MT-ART01 confinement, and does not assert Artin
holomorphy. The contrapositive of the just-cited local theorem
now proves that a possible pole must satisfy
\[
 \Re s_0=\tfrac12,\qquad q(s_0)\ge B(G)+1.
 \tag{MT-ART09a}
\]
If $4\nmid |G|$, use the strengthened bound at least four.
Thus, if the relative zeta function has order at most $B(G)$
at every point of the critical line, all Artin functions are
holomorphic throughout the open critical strip. This last
zero-order condition is an extra input, not supplied by $M_H$.
For an unconditional local conclusion at a particular $s_0$,
only the corresponding relative-order inequality is needed.

This is a strictly finer criterion than measuring the order
of $\zeta_L$ alone: by the definition of a meromorphic order,
\[
 q(s_0)=\operatorname{ord}_{s_0}\zeta_L
                 -\operatorname{ord}_{s_0}\zeta_K.
\]
The base contribution is removed exactly, including when the
base has a multiple zero. It is not necessary to posit any
noncoincidence of zeros. The argument is a complete input
transfer to the cited classical local theorem, not a new proof
of that theorem or of general Artin holomorphy.

Geometrically, $\zeta_L/\zeta_K$ is the $L$-function of the
augmentation Artin motive attached to the finite etale cover
after deleting ramified primes, with the standard inertia
factors restored. The relative virtual permutation class is
$[\mathbb Q[G]]-[\mathbb Q]$ in its representation $K_0$.
This belongs to the cote1 trace and SGA5 finite-correspondence
framework already fixed in AG-C1, and sharpens the
solvable part of Q-ARTIN-POS beyond MT-ART02. It does not turn a
virtual-character inequality into coefficientwise positivity.

